\cleardoublepage
\chapter{PENDAHULUAN}
\label{chap:pendahuluan}

\section{Latar Belakang}
Kegiatan campaign pada industri perbankan membutuhkan data yang dapat digunakan untuk menentukan target, menjalankan komunikasi, dan membaca hasil setelah campaign berjalan. Pada praktiknya, pekerjaan tersebut tidak berhenti pada pembuatan query. Data perlu dipahami berdasarkan konteks bisnis, diperiksa kualitasnya, diproses melalui alur yang dapat dijalankan ulang, lalu disajikan dalam bentuk yang dapat dibaca oleh unit terkait.

PT Bank Syariah Indonesia (Persero) Tbk atau BSI resmi beroperasi pada 1 Februari 2021 setelah penggabungan tiga bank syariah. Pada 2026, nama perseroan secara administratif menjadi PT Bank Syariah Indonesia (Persero) Tbk \autocite{bsi2026}. Selama Kerja Praktik di BSI, penulis ditempatkan pada fungsi yang berkaitan dengan Campaign \& Leads Management. Posisi ini membuat penulis berhadapan langsung dengan kebutuhan bisnis yang harus diterjemahkan menjadi kebutuhan data dan analisis.

Pekerjaan penulis berkembang dalam tiga arah. Pertama, penulis mengembangkan query SQL dan alur ETL untuk kebutuhan campaign. Kedua, penulis menyusun analisis deskriptif untuk membaca efektivitas campaign financing tahun 2026. Ketiga, penulis mengembangkan spatial clustering dari data log aplikasi digital untuk membaca pola lokasi pada tingkat agregat. Ketiga pekerjaan tersebut memiliki root yang sama, yaitu bagaimana data yang besar dan tersebar dapat diolah menjadi informasi yang dapat digunakan tanpa kehilangan konteks bisnis.

\section{Tujuan Kerja Praktik}
Tujuan Kerja Praktik ini adalah:
\begin{enumerate}
\item memahami alur kebutuhan data pada fungsi campaign dan leads;
\item menerapkan SQL dan ETL pada lingkungan big data perusahaan;
\item menyusun analisis deskriptif dari data campaign;
\item menerapkan spatial clustering pada data log aplikasi digital;
\item melatih validasi data, komunikasi kebutuhan bisnis, dan penerapan kerahasiaan data dalam pekerjaan.
\end{enumerate}

\section{Tempat dan Waktu Pelaksanaan}
Kerja Praktik dilaksanakan di PT Bank Syariah Indonesia (Persero) Tbk pada fungsi Campaign \& Leads Management. Lokasi kerja berada di BSI Tower, Jl. Medan Merdeka Selatan No. 17, RT 11/RW 2, Kelurahan Gambir, Kecamatan Gambir, Kota Jakarta Pusat, Daerah Khusus Ibukota Jakarta 10110.

Pelaksanaan Kerja Praktik berlangsung pada 6 Juli sampai 4 September 2026. Pembagian kegiatan utama selama periode tersebut dirangkum pada Tabel \ref{tab:jadwal-kp}. Tabel ini menunjukkan urutan pekerjaan, bukan catatan harian secara rinci.

\begin{table}[htbp]
\caption{Ringkasan jadwal pekerjaan selama Kerja Praktik}
\label{tab:jadwal-kp}
\centering
\small
\setlength{\tabcolsep}{6pt}
\renewcommand{\arraystretch}{1.30}
\begin{tabularx}{0.92\textwidth}{@{}
>{\raggedright\arraybackslash}X
>{\centering\arraybackslash}p{1.35cm}
>{\centering\arraybackslash}p{1.65cm}
>{\centering\arraybackslash}p{1.85cm}
@{}}
\toprule
\textbf{Kegiatan} & \textbf{Juli} & \textbf{Agustus} & \textbf{September} \\
\midrule
Pemahaman proses bisnis dan lingkungan data & X & X & \\
\addlinespace[2pt]
Pengembangan query SQL dan ETL & X & X & \\
\addlinespace[2pt]
Analisis deskriptif campaign financing & & X & X \\
\addlinespace[2pt]
Spatial clustering dan visualisasi & & X & X \\
\addlinespace[2pt]
Penyusunan laporan dan sanitasi informasi & & & X \\
\bottomrule
\end{tabularx}
\end{table}

\section{Ruang Lingkup dan Batasan}
Ruang lingkup laporan dibatasi pada pekerjaan yang benar-benar dilakukan penulis. Pekerjaan tersebut meliputi pengembangan query dan ETL, analisis deskriptif campaign financing, serta spatial clustering dan visualisasi. Social Network Analysis dibahas sebagai arah pengembangan karena pekerjaan tersebut belum selesai pada periode Kerja Praktik.

Laporan ini tidak mencantumkan data mentah, identifier nasabah, nomor telepon, nama tabel internal, nama kolom sensitif, query produksi secara lengkap, kode campaign, maupun angka bisnis internal yang bersifat terbatas. Nama produk atau unit hanya digunakan jika diperlukan untuk menjelaskan konteks secara umum. Pembatasan ini tidak mengubah metode yang dilaporkan, tetapi mengurangi detail yang dapat mengarah pada rekonstruksi data atau proses internal.

\section{Sistematika Penulisan}
Laporan mengikuti struktur template LaTeX terbaru yang digunakan pada mata kuliah II4090 Kerja Praktik. Bab I menjelaskan latar belakang, tujuan, tempat dan waktu, ruang lingkup, serta batasan. Bab II menjelaskan organisasi dan lingkungan kerja pada tingkat yang relevan terhadap pekerjaan penulis. Bab III menjelaskan pekerjaan yang dilakukan, mulai dari pemahaman data, pengembangan query dan ETL, analisis campaign, sampai spatial clustering. Bab IV berisi rangkuman pekerjaan, hasil, evaluasi diri, saran, dan ucapan terima kasih. Keterangan sanitasi informasi diletakkan pada lampiran agar isi utama tetap fokus pada pekerjaan.
